\begin{frame}[fragile]{실습: 대명사 지시 관계를 시각화}(self-attention)
  네 단어 $[\text{동물}, \text{도로}, \text{건너지},
  \text{그것은}]$ 에 직접 설계한 임베딩으로 self-attention:
\begin{lstlisting}
words = ["동물", "도로", "건너지", "그것은"]
E = [[3,0,0], [0,3,0], [0,0,3], [2.8,0.3,0]]
# "그것은"이 "동물"과 유사하게 설계됨
Q = K = V = E
output, weights = attention(Q, K, V, d_k=3)
for w, row in zip(words, weights):
    print(w, [round(x, 3) for x in row])
\end{lstlisting}
  \vskip0.2em
  \begin{columns}[T]
    \begin{column}{0.52\textwidth}
      \scriptsize
      \begin{tabular}{@{}lrrrr@{}}
        \toprule
        & 동물 & 도로 & 건너지 & 그것은 \\
        \midrule
        동물 & 0.582 & 0.003 & 0.003 & 0.412 \\
        도로 & 0.005 & 0.980 & 0.005 & 0.009 \\
        건너지 & 0.005 & 0.005 & 0.984 & 0.005 \\
        \textbf{그것은} & \textbf{0.561} & 0.007 & 0.004 & 0.427 \\
        \bottomrule
      \end{tabular}
    \end{column}
    \begin{column}{0.44\textwidth}
      \centering
      \includegraphics[width=\linewidth]{ch12_2_attention_heatmap.png}
    \end{column}
  \end{columns}
  {\footnotesize ``그것은'' 행이 ``동물'' 열(0.561)과 자기 자신
  열(0.427)에만 높은 값 — 대명사-선행사 관계가 숫자로 구현된
  모습. 나머지 행은 대각선(자기 자신)에 집중.}
\end{frame}
